\section{A determinant formula and two universal nonsplit-source barriers}
\label{sec:finite-source-barrier}

The splitness of the two preceding witnesses is forced by their local
geometry.  We prove a parabolic determinant formula and then show that
every dominant sixfold curve-cover family using either of these natural
representations has exactly the relevant involution inertia types.
The conclusion concerns these representations and does not classify all
finite groups with quadratic character field.

\subsection*{D-DET-01. The Seifert determinant formula}

Let \(K\) be an imaginary quadratic field, let \((V,h)\) be a positive
definite Hermitian \(K\)-space of dimension \(m\), and let
\(g_1,\ldots,g_r\in U(V,h)\) be finite-order elements with product one
and no common invariant vector.  We use right actions in this proof, to
match exactly the conventions of \cite[Theorem 2.5, p.~12]{G34:DWdual}.
Put
\[
 U_i=\operatorname{im}(g_i-1),\quad d_i=\dim_KU_i,\quad
 D=\sum_i d_i,\quad q_i=g_{i+1}\cdots g_r,
 \quad L_i=U_iq_i.
\]
Each restriction \(h_i=h|_{U_i}\) is nondegenerate.  Assume the
parabolic multiplicity space has even dimension \(2n=D-2m\).
Set
\[
                   b=\prod_i\det(1-g_i|_{U_i}).
\]
Then \(b\in\mathbb Q^*\), and the signed Hermitian discriminant of
the parabolic space is
\begin{equation}\label{eq:det-seifert}
       \boxed{\displaystyle
       \left[\frac{\prod_i\det h_i}{b}\right]
          \quad\text{in }\mathbb Q^*/N_{K/\mathbb Q}(K^*).}
\end{equation}
The determinants of the \(h_i\) may be computed in arbitrary
\(K\)-bases; changing those bases changes the expression by a norm.

\begin{proof}
Let \(\mathcal V=\bigoplus_iL_i\), with summation map
\(R:\mathcal V\to V\), \(R(w)=\sum_iw_i\).  A parabolic cocycle
\(u_i\in U_i\) has transported values \(w_i=u_iq_i\), and its relation
is \(R(w)=0\).  Its coboundary map is
\[
                    d(a)_i=a(g_i-1)q_i.
\]
The product-one relation makes \(Rd=0\).  The common-invariant-vector
hypothesis makes \(d\) injective.

Let \(G_i=q_i^{-1}g_iq_i\) act on \(L_i\), and set
\[
             A_i=(1-G_i^{-1})^{-1}|_{L_i}.
\]
These inverses exist, because \(L_i\) is the moving space.  Unitarity
gives \(A_i+A_i^*=1\).  On \(\mathcal V\) define the triangular
sesquilinear form
\[
 \mathcal S(w,z)=\sum_{j<i}h(w_j,z_i)+\sum_i h(w_iA_i,z_i).
\]
Its diagonal blocks are invertible, so it is nondegenerate.  Since
\(A_i+A_i^*=1\),
\[
                  \mathcal S+\mathcal S^*=R^*hR.
\]
We also have the crucial identity
\begin{equation}\label{eq:det-coboundary}
                     \mathcal S(d(a),z)=h(a,Rz).
\end{equation}
Indeed
\(\sum_{j<i}d(a)_j=a-a g_iq_i\).  Writing \(P_i\) for the
orthogonal projector onto \(U_i\), the diagonal term satisfies
\[
 d(a)_iA_i=a(g_i-1)(1-g_i^{-1})^{-1}q_i
                         =a g_iP_iq_i.
\]
Their sum differs from \(a\) by a vector in
\(\ker(g_i-1)q_i=L_i^\perp\).  Pairing with \(z_i\in L_i\) gives
\eqref{eq:det-coboundary}.

Because \(\mathcal S\) is nondegenerate and \(d\) is injective,
\eqref{eq:det-coboundary} also proves that \(R\) is surjective.
On \(\ker R\), \(\mathcal S\) is anti-Hermitian, with radical
exactly \(dV\).  Thus it induces a nondegenerate anti-Hermitian form
\(S_P\) on
\[
                         P=\ker R/dV.
\]

This is the geometric parabolic cup form up to one overall sign.
Here is the convention check, rather than an appeal to an unspecified
auxiliary form.  Apply Dettweiler--Wewers Theorem2.5 to the dual cocycle
obtained from \(h\).  Its formula simplifies, using
\(u_i'=u_i(g_i-1)^{-1}\), to
\[
 \sum_i h(u_i,v_i(g_i-1)^{-1})
       -\sum_{j<i}h(u_j,v_i(g_{j+1}\cdots g_i)^{-1}).
\]
Transporting both arguments by the tails \(q_i\), the second term is
\(-\sum_{j<i}h(w_j,z_i)\).  The first is
\(-\sum_i h(w_iA_i,z_i)\), by the adjoint identity for a unitary
operator.  Thus their cup pairing is \(-S_P\).  Since \(\dim_KP=2n\),
this sign has no effect on its determinant.  For involutions the diagonal
operator is exactly \(A_i=\tfrac12 I\), agreeing with the triangular
form used in the explicit certificates.

Choose a basis of \(\mathcal V\) consisting of a basis of \(dV\),
representatives of a basis of \(P\), and a complement \(C\) on which
\(R\) is an isomorphism.  In this basis the form has block shape
\[
 \begin{pmatrix}
 0&0&B\\0&S_P&F\\-B^*&-F^*&J
 \end{pmatrix},\qquad B=hR_C.
\]
Taking the Schur complement of \(S_P\) shows that its determinant is
\(N_{K/\mathbb Q}(\det B)\det S_P\).  A basis change of a
sesquilinear form multiplies its determinant by a norm.  Therefore
\([\det S_P]=[\det\mathcal S]\).

The triangular diagonal blocks give
\[
               \det\mathcal S=\frac{\prod_i\det h_i}{b}.
\]
For clarity, \(\det(h(\,\cdot\, A_i,\,\cdot\,))\)
equals \(\det h_i\,\overline{\det A_i}\), and unitarity makes
\(\overline{\det A_i}=\det(1-g_i|_{U_i})^{-1}\).
Moreover
\[
 \bar b=(-1)^D\left(\prod_i\det g_i\right)^{-1}b=b,
\]
because \(D=2m+2n\) is even and \(\prod_i g_i=1\).
Thus \(b\) and the determinant expression are rational.

Finally, multiply the anti-Hermitian form by any nonzero imaginary
\(\eta\in K\) to obtain a Hermitian form.  Its signed determinant is
\[
             (-1)^n\det(\eta S_P)
                   =N_{K/\mathbb Q}(\eta^n)\det S_P.
\]
It has the class claimed in \eqref{eq:det-seifert}.  The cup pairing on
the one-copy abelian factor differs from the multiplicity pairing only
by the rational cover-degree and root-line polarization scalars.  Such
a scalar to the power \(2n\) is a norm.  Switching the selected dual
or conjugate multiplicity convention also preserves the class (the
quotient \(\mathbb Q^*/N(K^*)\) has exponent two).  Hence the formula
applies to the actual polarized abelian factor.
\end{proof}

For involution inertia, \(b=2^D\) is a norm.  Thus
\begin{equation}\label{eq:det-involution}
           \boxed{[\operatorname{disc}_{\rm signed}P]
                    =\prod_i[\det(h|_{U_i})].}
\end{equation}
This formula uses only the local Hermitian moving spaces.  It is
independent of the ordering, Nielsen class, and Hurwitz component.

\subsection*{D-DET-02. The \(G_{24}\) barrier}

For the natural \(G_{24}\)-representation over
\(K=\mathbb Q(\sqrt{-7})\), every reflection has a norm-two root.
Since \(2=N_{K/\mathbb Q}(\tau)\), each reflection moving line has
norm determinant.  Thus every product-one generating tuple of these
involutions giving a balanced multiplicity space has split signed
discriminant.  This conclusion holds for unpaired tuples as well.

In fact every involution in this representation has norm moving-space
determinant.  Rank-one involutions are the reflections.  For a rank-two
involution, its negative plane is the orthogonal complement of the root
line of its negative, which is a reflection because \(-I\in G_{24}\).
Since \(\det H=1\), its plane determinant has class \(1/2\), again a
norm.  The remaining involution is \(-I\), whose moving space has
determinant one.  Thus \eqref{eq:det-involution} also covers arbitrary
involution types in \(G_{24}\).

There is a stronger full-source obstruction.  Let the base curve have
genus \(\gamma\), let \(r\) be the number of genuine branch points, and
suppose the natural multiplicity factor is a sixfold.  The parabolic
Euler characteristic gives
\[
             6=(2\gamma-2)\cdot3+D,\qquad D=12-6\gamma.
\]
Each nontrivial inertia element has moving-space dimension at least one,
so \(r\le D\).  For \(\gamma=1\), the Hurwitz dimension is at most
\(6\); for \(\gamma=2\), only unramified covers are possible and the
base moduli dimension is three; larger \(\gamma\) is impossible.
Thus a period map of rank nine must have \(\gamma=0\).

For \(\gamma=0\), the Hurwitz dimension is \(r-3\le9\).  Rank nine
forces \(r=12=D\), so all inertia elements have moving rank one.  Their
nontrivial eigenvalue belongs to \(K\), since it is the trace minus two.
The only roots of unity in \(K\) are \(\pm1\), so they are involutive
reflections.  The determinant formula forces splitness.

\textbf{Conclusion.} No dominant Weil-sixfold curve-cover family using
the natural \(G_{24}\)-representation can lie in a nonsplit determinant
component, even if arbitrary inertia orders and base genera were initially
allowed.  The inequalities are used only as necessary conditions for
dominance; positive rank nine was proved separately for the split witness.

\subsection*{D-DET-03. The \(F_4\)-extension barrier}

First we prove the local assertion for every outside involution, without
assuming it lies in a particular finite conjugacy class.  An element
\(g\in G_F\setminus W\) has the form \(tA\), with \(A\in M_4(\mathbb Q)\)
and \(t^2=-2\).  If \(g^2=1\), unitarity gives
\[
                       A^t=-A,\qquad A^2=-I/2.
\]
Let \(U=\ker(g+1)\).  The rational-linear map
\[
               x\longmapsto (1-tA)x:\mathbb Q^4\longrightarrow U
\]
is an isomorphism of rational vector spaces.  It lands in \(U\) by the
displayed square relation; injectivity follows by comparing rational and
\(t\)-coefficients.  Its Hermitian diagonal form is
\[
              h((1-tA)x,(1-tA)x)=2x^t x.
\]
Thus every such negative plane has the same underlying rational quadratic
form as the explicit negative plane of \(z\).

We use the following elementary trace-Witt fact.  If two \(K\)-Hermitian
forms have isometric rational diagonal forms, their difference has
hyperbolic rational diagonal form.  A nonzero rational isotropic vector
is also a \(K\)-isotropic vector.  Split off a Hermitian hyperbolic plane;
its underlying rational form is hyperbolic of dimension four.  Rational
Witt cancellation leaves a hyperbolic complement, and induction shows
the whole Hermitian difference is hyperbolic.  In particular equal-rank
forms have determinant classes differing by a norm.  This argument does
not identify a real signature with a rational Witt class.

The explicit plane of \(z\), with the basis in D-GF4-02, has Gram matrix
\[
                         \begin{pmatrix}2&t\\-t&2\end{pmatrix},
                         \qquad\det=2=N_{K/\mathbb Q}(t).
\]
Every outside involution therefore has norm moving-plane determinant.
The reflection moving lines have norms one or two, also norms from
\(K=\mathbb Q(\sqrt{-2})\).  Hence every product-one generating tuple
consisting of reflections and outside involutions has split discriminant,
whether or not its entries occur in adjacent equal pairs.

Now allow arbitrary inertia.  The sixfold dimension formula is
\[
                 6=(2\gamma-2)\cdot4+D,
                 \qquad D=14-8\gamma.
\]
For \(\gamma=1\), the Hurwitz dimension is at most six; for larger
\(\gamma\), the dimension equation is impossible.  Thus rank-nine
dominance requires a genus-zero base.

On that base, the product-one generating tuple must contain at least
two elements outside \(W\): the quotient \(G_F/W\) has order two,
so their number is positive and even.  For any outside element,
\(\bar g=-g\).  Complex conjugation therefore carries its \(+1\)
eigenspace to its \(-1\) eigenspace.  Its fixed dimension is at most
two, so its moving dimension is at least two.  Consequently
\[
                 r\le D-2=12.
\]
Dominance requires \(r\ge12\), so equality holds.  There are exactly
two outside elements, each with moving dimension two, and all ten other
elements have moving dimension one.  An outside element of moving rank
two has two eigenvalues \(+1\) and two \(-1\), so it is an involution.
The ten rank-one elements are reflections in \(W\).  The determinant
formula just proved therefore forces splitness.

\textbf{Conclusion.} No dominant Weil-sixfold curve-cover family using
this natural \(G_F\)-representation can lie in a nonsplit determinant
component, for any initially allowed inertia or base genus.

\subsection*{What a nonsplit construction would have to change}

For these two representations, changing the product-one tuple, Hurwitz
component, or ordering cannot produce a dominant nonsplit source.  Using
higher-order inertia or higher-genus bases within the same representations
loses the necessary nine-dimensional parameter space.  A new dominant
nonsplit source must therefore change the representation or the geometric
source construction.

Formula \eqref{eq:det-seifert} identifies the arithmetic alternative for
another representation: either a local moving-space determinant or the
product \(\prod_i\det(1-g_i|_{U_i})\) must contribute a nonnorm class.
For higher-order inertia, the automatic supplier argument also needs
replacement or a new check.  Cyclic quotients may then have local
exponents other than \(m/2\); their exponents need not occur in opposite
pairs, and their primitive embeddings need not all be balanced.  Schoen's
paired-exponent determinant construction and the full Brauer-character
coverage cannot simply be carried over without verifying those conditions.
This is an exact missing requirement, not a claim that all higher-order
suppliers fail.
